% §10.5 Bridge
%
% Cross-reference resolution status:
%
% DEFINED here:
%   \label{sec:ch10-bridge}                       -- this section
%
% RESOLVED (chapter 10, prior sections):
%   \Cref{sec:sketches-as-bernoulli}              -- §10.1
%   \Cref{sec:countmin}                           -- §10.2
%   \Cref{thm:countmin-bound}                     -- §10.2
%   \Cref{sec:hyperloglog}                        -- §10.3
%   \Cref{prop:hyperloglog-error}                 -- §10.3
%   \Cref{sec:minhash-lsh}                        -- §10.4
%   \Cref{prop:minhash-error}                     -- §10.4
%
% RESOLVED (chapter 9):
%   \Cref{ch:approximate-functions}               -- chapter 9
%   \Cref{def:n-ary-channel}                      -- §9.2
%   \Cref{thm:n-ary-kronecker}                    -- §9.3
%
% FORWARD REFS to later chapters (Plan 0 stubs):
%   \Cref{ch:randomized-algorithms}               -- chapter 11
%   \Cref{ch:approximate-relations}               -- chapter 12
%   \Cref{ch:algebra-types}                       -- chapter 13
%
% §10.5 is the short closing section of chapter 10. Three paragraphs:
% a recap of the three sketches as $\Bernoulli{T}$ instances, a forward
% pointer to chapter 13's composition algebra, and a forward pointer to
% chapters 11 and 12.

\section{Bridge}
\label{sec:ch10-bridge}

The three sketches built in this chapter share one skeleton and differ
only in the structure of the output type and the shape of the channel
matrix that ties true input to reported estimate. CountMin
(\Cref{sec:countmin}) is the frequency sketch: its output type is
$T = \mathbb{N}$, its query is one-sided (the estimate $\hat f_a$
never underestimates the true count $f_a$), and its channel matrix
$Q$ places all probability mass at or above the diagonal, with row
$f_a$ supported on $\{f_a, f_a + 1, \ldots\}$ and tail probability
controlled by $\Pr[\hat f_a \ge f_a + \epsilon F] \le \delta$ from
\Cref{thm:countmin-bound}. HyperLogLog (\Cref{sec:hyperloglog}) is the
cardinality sketch: its output type is again $T = \mathbb{N}$, but its
error is two-sided and approximately log-normal, with the standard
relative error $1.04 / \sqrt m$ of \Cref{prop:hyperloglog-error} fixing
the multiplicative width of the channel-matrix band around the
diagonal. MinHash (\Cref{sec:minhash-lsh}) is the similarity sketch:
its output type is the unit interval $T = [0, 1]$, its per-coordinate
agreement is a Bernoulli draw with success probability $J(A, B)$, and
the Hoeffding bound of \Cref{prop:minhash-error} controls the width of
the channel kernel around the diagonal at signature length $k$. Three
sketches, three $\Bernoulli{T}$ types, one common pattern: a
construction, an error analysis bounding $Q$ row-by-row, a composition
rule for combining sketches built on different streams.

The composition rules of the three sketches share a deeper structure
that the present chapter notes in passing but does not axiomatize.
Merging two CountMin tables is entry-wise integer addition on the
$d \times w$ counter array; merging two HyperLogLog register banks is
entry-wise maximum on the $m$-register array; merging two MinHash
signatures (in the natural sense of forming the signature of the union
set from the two component signatures) is entry-wise minimum on the
$k$-coordinate array. Each operation is associative, commutative, and
has an identity (the all-zero table, the all-zero register bank, the
all-$\infty$ signature respectively). Each operation respects the
chapter~9 channel structure: the merge of two $Q$-matrices is itself a
$Q$-matrix for the merged input, with the rows of the merged matrix
computable from the rows of the components by a known formula.
\Cref{ch:algebra-types} treats this pattern systematically. It
introduces the composition algebra of $\Bernoulli{T \times U}$
(product types, for sketches over pairs), $\Bernoulli{T + U}$ (sum
types, for sketches over tagged unions of input domains), and
$\Bernoulli{L\, T}$ (list types, for sketches whose output is a
sequence rather than a single value), and shows that the three
sketches of the present chapter are instances of those constructors.
The chapter~13 algebra is what makes the family of sketches
combinatorial rather than ad hoc: once the type constructors are
fixed, new sketches compose by structure.

Chapter~\ref{ch:randomized-algorithms} turns from streaming and
estimation to the algorithmic side of the same picture: randomized
algorithms whose answer is approximate in a one-sided sense, with
Miller-Rabin primality testing as the canonical worked example. The
move is from the data-structure setting of the present chapter (a
summary built from a stream, queried many times) to the
algorithm-execution setting (a Monte Carlo procedure run with fresh
randomness, possibly many times, against a single fixed input). The
channel framework of chapter~9 carries over without modification: a
randomized one-sided algorithm is a $\Bernoulli{\{\text{accept,
reject}\}}$ instance with a known false-positive or false-negative
rate, and amplification by repetition is composition of the same
channel matrix with itself. Chapter~\ref{ch:approximate-relations}
lifts the framework from approximate functions to approximate
relations, which is the necessary generalization for join queries and
multi-set membership tests. Chapter~\ref{ch:algebra-types}
unifies the three threads (the chapter~10 sketches, the chapter~11
randomized algorithms, the chapter~12 relations) under a single
algebraic vocabulary, and is the end of Part~III.
